\documentclass[12pt]{article}
\usepackage{comment}
\usepackage{soul}
\usepackage{graphicx}
\usepackage[dvipsnames]{xcolor}
\usepackage{amsmath,amsfonts,amssymb}
\usepackage{ulem}

\setlength{\textwidth}{18.2cm}
\setlength{\textheight}{21cm}
\setlength{\evensidemargin}{-1cm}
\setlength{\topmargin}{-1cm}
\setlength{\oddsidemargin}{-1cm}
\setlength\parindent{0pt}

\usepackage{color}
\definecolor{green2}{rgb}{0,0.56,0.32}

\newcommand{\lucien}[1]{\textcolor{cyan}{[Lucien: #1]}}
\newcommand{\cyan}[1]{\textcolor{cyan}{#1}}
\newcommand{\red}[1]{\textcolor{red}{#1}}
\newcommand{\scyan}[1]{\sout{\textcolor{cyan}{#1}}}
\newcommand{\Huayang}[1]{{\color{blue} [HS: #1]}}
\newcommand{\Felix}[1]{{\color{magenta} [FK: #1]}}
\newcommand{\jhyu}[1]{{\color{red} JH: #1}}
\newcommand{\Wei}[1]{{\color{green2} Wei: #1}}
\newcommand{\Shufang}[1]{{\color{red} SS: #1}}

\newcommand{\note}[1]{{\color{red} [Note: #1]}}


\newcommand{\changed}[2]{{\protect\color{red}\sout{#1}}{\protect\color{blue}\uwave{#2}}}


\begin{document}


%\title{Reply to Referee}
\begin{center}
\textsc{\bf \Large Reply to Report for LN19434}\\
\vspace{0.2in}
{\bf Title: Probing compressed Higgsinos at the FASER experiment}\\
\vspace{0.2in}
{\bf Authors:  Shufang Su, Wei Su, Jin Min Yang, Pengxuan Zhu, Rui Zhu
}
\end{center}

 
We would like to thank the referees for careful reading of the manuscript and for the constructive suggestions.   We have carefully considered the questions/comments from the referee that certainly have helped improve the manuscript. Below are our responses and revisions. \\

%\smallskip

%$\bullet$
%\begin{numerate}

\noindent

% ------------------------------------------
{\it 1. `` My main concern is that the central
claim relies on a quantitatively delicate prediction of an extremely small neutral Higgsino mass splitting, while the analysis does not treat this ingredient with the level of precision required by the claimed signal region.

\quad The key weakness of the manuscript is the treatment of the neutral Higgsino mass difference. The phenomenologically interesting region identified by the authors lies in the range $\Delta m_0 \sim 4-30$ MeV, but the analysis is based primarily on approximate tree-level expressions together with the qualitative statement that one-loop corrections mainly shift the cancellation condition without altering the overall conclusion. This is not sufficient for the present problem. The signal region is not simply a broadly compressed electroweakino spectrum; rather, it is a narrow cancellation-driven region in which the observable quantity is the small residual difference left after cancellation among much larger contributions. In such a situation, even modest loop corrections can shift $\Delta m_0$ by an amount comparable to, or larger than, the quoted signal window. Since the decay width is highly sensitive to the mass difference, the lifetime and therefore the predicted FASER event yield may change substantially even under a small shift in the splitting.

\quad For this reason, I do not find the manuscript’s treatment of loop effects adequate. In my view, the issue is not merely that the authors have not presented a full precision calculation;
rather, the more basic problem is that, in such an extreme near-degenerate regime, even a
nominal one-loop treatment may well be insufficient unless accompanied by a convincing assessment of missing higher-order effects and of the numerical stability of the spectrum itself. The manuscript assumes, without demonstrating it, that loop corrections merely shift the cancellation locus without materially affecting the phenomenologically relevant 4–30 MeV window. In such a cancellation-driven regime this assumption is not justified. Moreover, the accuracy of standard spectrum solvers in resolving splittings at the level of only a few MeV is itself open to question. A quantitatively reliable analysis would therefore require a substantially more careful treatment of the electroweakino spectrum, including at least a clearly defined loop-level calculation, a discussion of the dependence on renormalization prescription, a demonstration of numerical stability in the cancellation region, and a credible estimate of the uncertainty from uncomputed higher-order effects. Without such a treatment, the central reach claim does not appear sufficiently robust." } \\
%{\bf Higher energy lepton colliders}

REPLY:   \\

\smallskip
\noindent

% ------------------------------------------
{\it 2. `` A related concern is that the phenomenologically relevant region appears to arise only for a rather special choice of MSSM parameters, such that a strong cancellation occurs in the neutral Higgsino splitting. This means that the signal region corresponds to a narrow and highly tuned slice of parameter space rather than to a generic compressed Higgsino scenario. A tuned region is not automatically uninteresting, but if the authors wish to make a strong case for the broader importance of this scenario, they should either quantify the degree of tuning explicitly or provide a concrete motivation from a high-scale or SUSY-breaking perspective for why the relevant gaugino-mass relation should arise. At present, the manuscript does not do so, which weakens the broader phenomenological significance of the proposal. " } \\

REPLY:   \\
\smallskip
\noindent

% ------------------------------------------
{\it 3. `` Independently of the issue above, I find the detector-level treatment of the proposed signal
insufficient. The visible signature is the photon from $\chi_2^0 \to \chi_1^0 \gamma$. In the rest frame of $\chi_2^0 $, the photon energy is approximately set by $\Delta m_0$, which in the claimed signal region is only 4–30 MeV. The actual observability of such a signal at FASER/FASER-2 depends on several nontrivial detector-level issues, including the laboratory-frame photon energy spectrum, the angular distribution of the photon, the probability of photon conversion, trigger ability, photon reconstruction efficiency, and possible geometric losses. The manuscript demonstrates acceptance of the long-lived neutralino into the detector volume, but not the acceptance and reconstruction of the visible decay product that is supposed to establish the signal. This is a significant omission." } \\

REPLY:  

\smallskip
\noindent


% ------------------------------------------
{\it 4. ``A further weakness is the absence of a serious treatment of backgrounds. The projected sensitivity appears to be based essentially on geometric acceptance, decay probability inside the fiducial volume, and an assumed three-event reach. This may be acceptable as a first exploratory estimate, but it is not sufficient for a strong phenomenological claim at the level targeted by the manuscript. For a mono-photon displaced signature, a realistic analysis should address at least the main potential background classes, such as neutrino-induced interactions, muon-induced secondaries or bremsstrahlung, and photon reconstruction or trigger limitations in the relevant soft-photon regime. In addition, the projected reach should be based on a detector-level signal definition with a quantified background model and a clear statistical procedure, rather than on an implicit three-event criterion. As it stands, the manuscript does not establish that the proposed mono-photon signature is experimentally realizable at FASER-2 with the claimed sensitivity, because neither the dominant backgrounds nor the photon efficiency in the relevant energy regime are quantitatively demonstrated." } \\

REPLY:  

\smallskip
\noindent


% ------------------------------------------
{\it 5. ``" } \\

REPLY:  

\smallskip
\noindent
 We thank the referee again for the careful reading. With those revisions, we feel that we have addressed  the referee's comments ...... \\


{\bf References }  \\

% [1] 

[1] 

\end{document}







